\documentclass[aspectratio=169,11pt]{beamer}
\makeatletter
\input{../../templates/dune-professional/beamerthemeDUNEProfessional.sty}
\makeatother
\title{VD detector protection}
\subtitle{Extending the BDE PLC system}
\author{Manuel Arroyave}
\date{I\&I workshop draft --- 27 October 2026}
\newcommand{\mydate}{27 October 2026}
\hypersetup{colorlinks=true,urlcolor=DuneInfo,linkcolor=DuneInfo}
\newcommand{\planningnote}[1]{\par\medskip{\footnotesize\color{DuneTextMuted}#1\par}}
\begin{document}

\begin{frame}[plain]
  \titlepage
\end{frame}

\begin{frame}{BDE already has the PLC and panel design}
  \begin{columns}[T,onlytextwidth]
    \begin{column}{0.43\textwidth}
      \colorbox{DunePaper}{\includegraphics[page=2,width=\dimexpr\linewidth-2\fboxsep\relax,height=3.35cm,keepaspectratio,
        trim=290bp 500bp 900bp 430bp,clip]{figures/tpc-panel-reference.pdf}}\par
      \medskip\DuneDiagramCaption{TPC protection panel drawings.}
    \end{column}
    \begin{column}{0.53\textwidth}
      Beckhoff PLC, distributed EtherCAT I/O and hardware supply inhibits.\par
      \bigskip
      VD allocation: 10 BDE auxiliary racks and 20 PL506 supplies for 80 WIECs.\par
      \bigskip
      Panel drawings, parts lists, power calculations and the Ignition front end already exist.
    \end{column}
  \end{columns}\par
  \medskip
  \textbf{Proposal: extend this design to VD, using industrial parts we can maintain and expand.}
  \planningnote{BDE = bottom drift electronics; PLC = programmable logic controller; WIEC = warm interface electronics crate. Sources: TPC panel package, VD rack workbook and June 2026 SC presentations.}
\end{frame}

\begin{frame}{Put the VD PLC in the CUC, with I/O near the equipment}
  \centering
  \begin{tikzpicture}[
    box/.style={fill=DunePaper,text=DuneOnAccent,align=center,
      text width=2.7cm,minimum height=1.05cm,inner sep=6pt},
    edge/.style={-{Stealth},draw=DuneSky,line width=1.2pt}]
    \node[box] (sc) at (-4.4,1.3) {Slow controls\\Ignition};
    \node[box,fill=DuneSky] (plc) at (0,1.3) {VD DPS PLC\\CUC proposed};
    \node[box] (facility) at (4.4,1.3) {Facility / cryogenic\\conditions};
    \node[box] (bde) at (-3.3,-1.1) {BDE and bias HV\\existing I/O design};
    \node[box,fill=DuneCoral] (other) at (3.3,-1.1) {PDS, TDE, CalCI\\additional interfaces};
    \draw[edge,<->] (sc) -- node[above,font=\small] {OPC UA} (plc);
    \draw[edge,dashed] (facility) -- (plc);
    \draw[edge,<->] (plc.south west) -- (-1,-0.1) -| (bde.north);
    \draw[edge,<->] (plc.south east) -- (1,-0.1) -| (other.north);
  \end{tikzpicture}\par
  \medskip
  \DuneDiagramCaption{Proposal: place the central PLC beside the other PLC racks in the CUC; retain distributed I/O.}
  \planningnote{CUC = Central Utility Cavern. Dashed link: facility signals need a selected hardware or software path. PDS = photon detection; TDE = top drift electronics; CalCI = calibration and cryogenic instrumentation. HD can use an extension or a separate PLC later.}
\end{frame}

\begin{frame}{Add the VD interfaces that the BDE system does not cover}
  \renewcommand{\arraystretch}{1.3}
  \begin{tabularx}{\textwidth}{@{}>{\RaggedRight\arraybackslash}p{3.2cm}>{\RaggedRight\arraybackslash}X@{}}
    \toprule
    \textbf{Known VD scope} & \textbf{Proposed DPS addition}\\
    \midrule
    PDS: 80 PoF boxes; 58 LV feeds & Local PoF readout/control; supply inhibits and trip feedback\\
    TDE: 20 racks & Crate-status exchange and external inhibit interface\\
    CalCI & Level, temperature and optical-source status; operation permits\\
    BMS / cryogenics & Selected facility conditions used by DPS; monitoring shared with SC\\
    \bottomrule
  \end{tabularx}\par
  \medskip
  The 80 PoF boxes produce 1,280 monitoring values. Local readout handles that
  volume; the protection permits determine the DPS I/O.\par
  \planningnote{Working VD inventory and ICD matrix; exact trip scope remains subsystem-specific. BMS = building management system. SDSTA and LBNF/DUNE have separate BMSs.}
\end{frame}

\begin{frame}{Protection must keep working when SC is unavailable}
  \begin{columns}[T,onlytextwidth]
    \begin{column}{0.48\textwidth}
      \DuneFixedColorBlock{paper}{2.65cm}{SC}{Operate and monitor}
        {Commands, sequencing, displays and history. An operating request cannot override a protection trip.}
    \end{column}
    \begin{column}{0.48\textwidth}
      \DuneFixedColorBlock{sky}{2.65cm}{DPS / local protection}{Protect equipment}
        {PLC permits and device protection remain effective through an SC outage. Reset and re-enable are separate actions.}
    \end{column}
  \end{columns}\par
  \bigskip
  The PX3 setup already has local fallback. The VD prototype will exercise it together with PLC protection and SC recovery.
  \planningnote{Proposed control model. The PX3 guide uses Ignition for normal protection and PDU fallback when Ignition is unreachable; this is a concrete integration test with Geoff and the rack team.}
\end{frame}

\begin{frame}{Test the full protection path before buying the additions}
  \centering
  \begin{tikzpicture}[
    box/.style={fill=DunePaper,text=DuneOnAccent,align=center,
      text width=2.25cm,minimum height=1.0cm,inner sep=5pt},
    edge/.style={-{Stealth},draw=DuneSky,line width=1.2pt}]
    \node[box] (in) at (-4.8,0) {Sensor or equipment\\fault input};
    \node[box,fill=DuneSky] (plc) at (-1.6,0) {PLC and\\remote I/O};
    \node[box,fill=DuneCoral] (out) at (1.6,0) {Supply inhibit\\or trip output};
    \node[box] (state) at (4.8,0) {Equipment state\\and feedback};
    \draw[edge] (in) -- (plc);
    \draw[edge] (plc) -- (out);
    \draw[edge] (out) -- (state);
  \end{tikzpicture}\par
  \medskip
  \DuneDiagramCaption{Use the BDE vertical-slice work and G271 equipment for representative VD paths.}
  \bigskip
  Measure detection through completed shutdown against the equipment's allowed
  time. Exercise SC, fieldbus and power loss, then reset and restart.\par
  \bigskip
  No critical sub-ms function has been identified so far. The simulation prepares load and fault cases; the bench measures hardware behaviour.
\end{frame}

\begin{frame}{Take the VD hardware to review in June 2027}
  \renewcommand{\arraystretch}{1.4}
  \begin{tabularx}{\textwidth}{@{}>{\RaggedRight\arraybackslash}p{3.2cm}>{\RaggedRight\arraybackslash}X@{}}
    \toprule
    \textbf{Work before review} & \textbf{Result for procurement}\\
    \midrule
    Adapt the BDE design & VD panels, added I/O, CUC fieldbus routes and power allocation\\
    Test the interfaces & Trip-time, failure and recovery results; SC/PDU integration\\
    Prepare purchases & Drawings, parts and cable lists, costs, lead times and spares\\
    \bottomrule
  \end{tabularx}\par
  \bigskip
  After readiness review: order the hardware and install it in the site sequence.
  \planningnote{Possible facility-interface need: summer 2027. Full VD response-matrix review: early 2028. The early matrix work supplies requirements that affect hardware purchases.}
\end{frame}

\begin{frame}{I will lead the VD hardware work with a control engineer}
  \renewcommand{\arraystretch}{1.4}
  \begin{tabularx}{\textwidth}{@{}>{\RaggedRight\arraybackslash}p{4.1cm}>{\RaggedRight\arraybackslash}X@{}}
    \toprule
    \textbf{Proposed contribution} & \textbf{Work}\\
    \midrule
    Manuel + control engineer & VD design, SURF interfaces and procurement package\\
    Terri and Trevor & BDE panels, PLC/I/O implementation and bench work\\
    Geoff + Manuel & SC commands, PDU protection and recovery\\
    Vivian + consortia / DPS & Equipment requirements and ICD updates\\
    \bottomrule
  \end{tabularx}\par
  \bigskip
  \textbf{Support needed:} control-engineering time and access to the existing test stands.
  \planningnote{SURF work: inspection maps are ready; access is being coordinated with Rachel, with Dale Curran identified as an SDSTA integration contact.}
\end{frame}

\appendix
\begin{frame}{Backup: VD rack space already allocated to DPS}
  \renewcommand{\arraystretch}{1.4}
  \begin{tabularx}{\textwidth}{@{}>{\RaggedRight\arraybackslash}p{3.4cm}r>{\RaggedRight\arraybackslash}X@{}}
    \toprule
    \textbf{Rack type} & \textbf{Count} & \textbf{DPS equipment / connections}\\
    \midrule
    BDE auxiliary & 10 & 12U per rack; eight WIEC connections per rack\\
    BDE low voltage & 20 & One PL506 per rack; four WIECs per supply\\
    Bias high voltage & 10 & 12U per rack for the DPS interface to ISEG supplies\\
    \bottomrule
  \end{tabularx}\par
  \bigskip
  The auxiliary racks contain the WIEC I/O; the bias racks contain the ISEG inhibit outputs.
  \planningnote{May 2025 VD rack workbook; June 2026 BDE layout. U is rack height. These are design allocations; the bias-HV sheet is a first estimate.}
\end{frame}

\begin{frame}{One BDE I/O group serves eight WIECs}
  \centering
  \begin{tikzpicture}[
    box/.style={fill=DunePaper,text=DuneOnAccent,align=center,
      text width=3.1cm,minimum height=1.05cm,inner sep=6pt},
    edge/.style={-{Stealth},draw=DuneSky,line width=1.2pt}]
    \node[box,fill=DuneSky] (io) at (-4,0) {BDE auxiliary rack\\remote I/O};
    \node[box] (lv) at (1.6,1.1) {2 LV racks\\2 PL506 supplies};
    \node[box,fill=DuneCoral] (wiec) at (1.6,-1.1) {8 WIECs\\8 PTCs + 48 WIBs};
    \draw[edge] (io.north east) -| (-1.7,1.1) -- node[above,font=\small] {inhibits} (lv.west);
    \draw[edge,<->] (io.south east) -| (-1.7,-1.1) -- node[below,font=\small] {EtherCAT} (wiec.west);
    \draw[edge] (lv.east) -- (4.35,1.1) -- node[right,font=\small] {48 V} (4.35,-1.1) -- (wiec.east);
  \end{tikzpicture}\par
  \medskip
  \DuneDiagramCaption{Ten groups: 20 PL506 supplies, 80 WIECs and 480 WIBs.}
  \medskip
  \raggedright The auxiliary rack also connects the fans, heaters and temperature sensors.
  \planningnote{WIEC = warm interface electronics crate; PTC = power and timing card; WIB = warm interface board. June 2026 layout; fan/heater wiring omitted here.}
\end{frame}

\begin{frame}{Panel drawings and the Ignition front end are available}
  \begin{columns}[T,onlytextwidth]
    \begin{column}{0.47\textwidth}
      \textbf{TPC protection panel designs}\par\medskip
      \colorbox{DunePaper}{\includegraphics[page=2,width=\dimexpr\linewidth-2\fboxsep\relax,height=3.1cm,keepaspectratio,
        trim=290bp 500bp 900bp 430bp,clip]{figures/tpc-panel-reference.pdf}}\par
      \medskip\DuneDiagramCaption{Central PLC, bias I/O and LV/fan/heater panels.}
    \end{column}
    \begin{column}{0.49\textwidth}
      \textbf{BDE Ignition screens}\par\medskip
      \includegraphics[page=7,width=\linewidth,height=3.1cm,keepaspectratio,
        trim=10bp 210bp 460bp 74bp,clip]{figures/sc-progress-june2026.pdf}\par
      \planningnote{Existing display of WIEC groups; June 2026 SC progress presentation.}
    \end{column}
  \end{columns}\par
  \medskip
  Parts list: CX5140 PLC, EK1100 couplers, EK1521 fibre links and EL2124 inhibit outputs.
  \planningnote{TPC reference: EDMS 2401090 v4, sheet 2; component names from the 2023 parts list.}
\end{frame}

\begin{frame}{The VD bias-panel draft specifies two DIN rails per rack}
  \renewcommand{\arraystretch}{1.3}
  \begin{tabularx}{\textwidth}{@{}>{\RaggedRight\arraybackslash}Xrr@{}}
    \toprule
    \textbf{Component} & \textbf{Per rail} & \textbf{10 racks}\\
    \midrule
    EK1100 EtherCAT coupler & 1 & 20\\
    EL9410 supply terminal & 1 & 20\\
    EL2124: four 5 V outputs & 18 & 360\\
    EL9011 end cap & 1 & 20\\
    24 V power supply & 1 & 20\\
    \bottomrule
  \end{tabularx}\par
  \medskip
  Each rail occupies 337 mm. The draft provides 1,440 output channels;
  the working VD inventory lists 576 bias channels.\par
  \planningnote{May 2025 workbook, Bias HV Rack!K17:K21. Quantities describe the first-pass panel layout.}
\end{frame}

\begin{frame}{VD power estimates for the BDE racks}
  \renewcommand{\arraystretch}{1.4}
  \begin{tabularx}{\textwidth}{@{}>{\RaggedRight\arraybackslash}p{3.7cm}>{\RaggedRight\arraybackslash}p{2.4cm}>{\RaggedRight\arraybackslash}X@{}}
    \toprule
    \textbf{Equipment} & \textbf{Per rack} & \textbf{VD basis}\\
    \midrule
    PL506 input & 2.25 kW & 20 racks: about 45.1 kW for the WIEC supplies\\
    Auxiliary DPS input & About 1.11 kW & 10 racks: about 11.1 kW; eight WIECs per rack\\
    Bias DPS supply allowance & 24 V, 40 A & 960 W DC capacity per rack in the workbook\\
    \bottomrule
  \end{tabularx}\par
  \bigskip
  The auxiliary estimate scales the recorded 831 W from six to eight WIECs:
  $831\times8/6 = 1{,}108$ W.\par
  \planningnote{Planning estimates from the May 2025 workbook. Linear scaling follows its VD note; it is not a measured load. Supply capacity is separate from AC demand.}
\end{frame}

\begin{frame}{The PX3 rack-protection setup already defines shutdown}
  \renewcommand{\arraystretch}{1.45}
  \begin{tabularx}{\textwidth}{@{}>{\RaggedRight\arraybackslash}p{3.7cm}>{\RaggedRight\arraybackslash}X@{}}
    \toprule
    \textbf{June 2026 setup guide} & \textbf{Documented behaviour}\\
    \midrule
    Ignition connected & Ignition reads sensors and switches outlets over SNMP\\
    Ignition unreachable & PX3 local rules handle smoke, water and temperature; the network-switch outlet stays on\\
    After a fallback trip & Power is restored manually or through Ignition\\
    \bottomrule
  \end{tabularx}\par
  \bigskip
  Test this setup together with the PLC's protection permits and trip feedback.
  \planningnote{DUNE-doc-32655 v8, pp.~3--6 and 14--16. The PDU fallback is documented; the PLC permit integration is the proposed extension.}
\end{frame}

\begin{frame}{The VD inventory covers the other subsystem connections}
  \renewcommand{\arraystretch}{1.35}
  \begin{tabularx}{\textwidth}{@{}>{\RaggedRight\arraybackslash}p{3.2cm}>{\RaggedRight\arraybackslash}X@{}}
    \toprule
    \textbf{Subsystem} & \textbf{Working inventory and interfaces}\\
    \midrule
    PDS readout / calibration & 42 DAPHNE boards and 8 calibration boxes; 58 LV feeds from MPOD modules\\
    Power over fibre (PoF) & 80 laser boxes; 640 temperature and 640 voltage readings; compact PLC interface proposed\\
    TDE & 20 mezzanine racks; power supplies and readout-crate management\\
    Calibration / cryogenic instrumentation & Level and temperature readout, purity monitors, ionization laser and neutron source\\
    \bottomrule
  \end{tabularx}
  \planningnote{VD hardware lists and SC architecture. Proposed routing: monitoring to Ignition through native drivers and OPC UA; protection permits and fault status exchanged with the local equipment.}
\end{frame}

\begin{frame}{The BMS drawing identifies 20 control-panel symbols}
  \renewcommand{\arraystretch}{1.4}
  \begin{tabularx}{\textwidth}{@{}>{\RaggedRight\arraybackslash}p{3.4cm}>{\RaggedRight\arraybackslash}X@{}}
    \toprule
    \textbf{Drawing area} & \textbf{Interfaces on the inspection maps}\\
    \midrule
    CUC, level 4850 & Eight panel references: air handling, chilled water, power and associated equipment\\
    Level 4910 & North and south sump interfaces, SP-1 and SP-2\\
    4850-08-ES group & Spare panel for future equipment; mapped to the access-drift inspection area\\
    Other 4850 areas & Chillers, pumps, ventilation and generator interfaces\\
    \bottomrule
  \end{tabularx}\par
  \medskip
  Inspection planned with Rachel; Dale Curran is the SDSTA integration contact identified for this work.
  \planningnote{U-FD-M-901 rev12, 14 February 2025; marked 4850/4910 maps and survey register. Symbols are drawing references; field observations are pending.}
\end{frame}

\begin{frame}{Reference: the May 2025 VD rack inventory}
  \renewcommand{\arraystretch}{1.15}
  \begin{tabularx}{\textwidth}{@{}>{\RaggedRight\arraybackslash}Xr>{\RaggedRight\arraybackslash}Xr@{}}
    \toprule
    \textbf{Rack type} & \textbf{Count} & \textbf{Rack type} & \textbf{Count}\\
    \midrule
    BDE low voltage & 20 & PDS low voltage & 4\\
    BDE auxiliary & 10 & Drift HV & 2\\
    TDE & 20 & Reserved drift HV & 2\\
    Bias HV & 10 & Spare / CalCI & 8\\
    \bottomrule
  \end{tabularx}\par
  \bigskip
  \textbf{76 numbered mezzanine positions:} SEDM-01 to SEDM-76.\par
  \medskip
  The worksheet also has a separate concept for 40 top-plate mini-racks.
  \planningnote{May 2025 allocation: rack positions, distinct from PLC panels and remote-I/O stations. Includes the listed reserved and spare positions.}
\end{frame}

\end{document}
